\documentclass[10pt,a4paper]{amsart}
\usepackage[margin=19mm]{geometry}
\usepackage{amsmath,amssymb,mathtools}
\usepackage[scr=rsfs,cal=euler]{mathalfa}
\usepackage{xfrac}
\usepackage[unicode,hidelinks,bookmarks=false]{hyperref}
\newcommand{\ie}{i.e.\ }
\newcommand{\cf}{cf.\ }
\newcommand{\resp}{respectively\ }
\newcommand{\Subsection}{\subsection}
\providecommand{\nobreakdash}{\nobreak\hbox{-}}
\setlength{\emergencystretch}{2em}
\multlinegap0pt
\usepackage{tikz}
\usepackage{amssymb}
\usepackage{xcolor}
\usepackage{dsfont}
\usepackage[shortlabels]{enumitem}
\newtheorem{lemma}{Lemma}
\newcommand{\ma}{\mathbf{a}}
\title{The spectral characteristics of the Sturm Hamiltonian with eventually periodic type}
\author{Jie CAO, Zhenyu Yu$^*$}
\date{Source: arXiv:2505.00328v2 (CC BY 4.0)}
\begin{document}
\maketitle

\noindent\emph{Source and licence.} The spectral characteristics of the Sturm Hamiltonian with eventually periodic type. Version arXiv:2505.00328v2 (CC BY 4.0), distributed by arXiv under the Creative Commons Attribution 4.0 International licence, http://creativecommons.org/licenses/by/4.0/, as recorded in the arXiv licence metadata for this exact version and shown on the abstract page https://arxiv.org/abs/2505.00328v2. Full licence notice: \texttt{licenses/arxiv-2505.00328-CC-BY-4.0.txt}.\par\smallskip

\noindent\textit{Editorial scope.} Counted unit: the article's lemma primitive (source0.tex lines 1261--1263). The article's complete original proof of it is delivered with it (source0.tex lines 1264--1329, one \texttt{proof} environment of 66 lines). No mathematics is written, repaired or re-derived: the statement and the proof are the article's own lines, in the article's own notation, and no retained line is substituted or re-flowed.  Retained uncounted as the source-local context the proof needs: lines 795--800.  Nothing was omitted from any retained span, so the package carries no omission record.  No retained line contains a citation, so the wrapper carries no bibliography and no bibliography entry.  No retained line contains a figure environment, an image-inclusion command or drawing code, so no image is delivered and the package declares no asset dependency.  Editorial formatting only, in addition: an AMS \texttt{amsart} preamble that restates the article's own declarations and macros used by the retained lines, namely the environments \texttt{lemma} on separate counters, together with the standard packages those lines need.  The retained environments print in this document as Lemma 1 in order of appearance, whereas the article prints the same items with the numbering of its own full document.  The complete native source of the article as distributed in the arXiv e-print for this exact version is bundled unchanged as \texttt{sources/177397/source0.tex}.
	\begin{align}\label{q-n}
		\begin{cases}
			p_{-1}=1, p_0=0,\  &p_{n+1}=a_{n+1}p_n+p_{n-1},\ n\geq0, \\
			q_{-1}=0, q_0=1,\ &q_{n+1}=a_{n+1}q_n+q_{n-1},\ n\geq0.
		\end{cases}
	\end{align}

	\begin{lemma}\label{primitive}
		$\hat{A}_{\ma},A_{\ma}$ are primitive and have the same Perron-Frobenius eigenvalue $E_{\ma}$. Moreover, we have $q_{kn}(\alpha)\sim E_{\ma}^n$ for all $\alpha\in\mathcal{P}(\ma)$, 
	\end{lemma}

	\begin{proof}
		It is straightforward to check that all the entries of $\hat{A}_{\ma}^5$ are positive,
		thus $\hat{A}_{\ma}$ is primitive. 
		
		Note that there exist matrices $$
		P=\left(\begin{array}{ccc}
			1&0&1\\
			-1&1&0\\ 
			-1&1&-1
		\end{array}\right)\ \ \text{and}\ \ Q_m=\left(%
		\begin{array}{cc}
			m & 1 \\ 
			1 & 0%
		\end{array}%
		\right)$$
		such that 
		\begin{equation*}
			\hat{A}_{\ma}=P\left(\begin{array}{cc}
				(-1)^k&0\\
				0&B_{\ma} 
			\end{array}\right)P^{-1},\ \  \text{where}\ B_{\mathbf{a}}=Q_{a_k}\cdots Q_{a_{2}}Q_{a_1}.
		\end{equation*}
		It follows that 
		$$\text{det}(\lambda I_3-\hat{A}_{\ma})=(\lambda-(-1)^k)\text{det}(\lambda I_2-B_{\ma}).$$
		Thus $\hat{A}_{\ma}$ and $B_{\ma}$ have the same Perron-Frobenius eigenvalue $E_{\ma}$. On the other hand, if we consider the graph related to the incidence matrix $A_{\ma}$, then it is easy to show that the graph is aperiodic. Consequently $A_{\ma}$ is primitive. Let $N=\#\mathscr{A}_{\ma}$, by direct computation,
		$$\text{det}(\lambda I_N-A_{\ma})=\lambda^{N-3}\text{det}(\lambda I_3-\hat{A}_{\ma}).$$
		Thus the Perron-Frobenius eigenvalue of $A_{\ma}$ is also $E_{\ma}$.
		
		Fix $\alpha\in\mathcal{P}(\ma)$ and by \eqref{q-n}, for any $n\in{\mathbb{N}%
		}$, we have 
		\begin{equation*}
			\left( 
			\begin{array}{c}
				q_{kn}(\alpha) \\ 
				q_{kn-1}(\alpha)
			\end{array}
			\right)=\left(\left( 
			\begin{array}{cc}
				a_k & 1 \\ 
				1 & 0%
			\end{array}
			\right)\cdots\left( 
			\begin{array}{cc}
				a_1 & 1 \\ 
				1 & 0%
			\end{array}
			\right)\right)^n\left( 
			\begin{array}{c}
				1 \\ 
				0%
			\end{array}
			\right)=B_{\ma}^n\left( 
			\begin{array}{c}
				1 \\ 
				0%
			\end{array}
			\right).
		\end{equation*}
		From this it is easy to show that there exist two constants $c_{\mathbf{a}%
		},d_{\mathbf{a}}$ such that 
		\begin{equation*}
			q_{kn}(\alpha)=c_{\mathbf{a}}E_{\mathbf{a}}^n+d_{\mathbf{a}}(-E_{\mathbf{a}%
			})^{-n}\sim E_{\mathbf{a}}^{n}. \label{qE}
		\end{equation*}
		The proof is complete.
	\end{proof}

\end{document}
